% Pregunta 610 (2009-octubre-OEP2009, n.º 7): Leontief y = min(b1 x1, b2 x2).
% Isocuantas en L con el vértice sobre el rayo x2/x1 = b1/b2 (aquí 0,8). Dos isocostes de
% pendientes -w1/w2 distintas (-2 y -0,35) tocan la isocuanta y0 en el mismo vértice: la
% combinación de factores no cambia con los precios (opción a).
% Autor: Víctor Gutiérrez Marcos
\documentclass[tikz,border=6pt]{standalone}
\input{tikz-tcee}
\begin{document}
\begin{tikzpicture}[tcee]
  \ejes{$x_1$}{$x_2$}{6.4}{5.2}
  % Rayo de los vértices
  \draw[guia] (0,0) -- (5.9,4.72) node[etiqueta,above,text=tceegris] {$\dfrac{x_2}{x_1}=\dfrac{\beta_1}{\beta_2}$};
  % Isocuantas en L
  \draw[curva1] (2,4.8) -- (2,1.6) -- (6,1.6) node[pos=0.55,etiqueta,above,text=tceemorado] {$y_0$};
  \draw[curva1] (3.5,4.8) -- (3.5,2.8) -- (6,2.8) node[etiqueta,above left,text=tceemorado] {$y_1>y_0$};
  % Isocostes por el mismo vértice
  \draw[curva2] (2.8,0) -- (0.65,4.3) node[etiqueta,above,text=tceeazul] {$w_1/w_2$ alto};
  \draw[curvanueva] (0,2.3) -- (5.9,0.235) node[pos=0.78,etiqueta,above=4pt,text=tceerojo] {$w_1/w_2$ bajo};
  % Combinación elegida
  \proyeccion{(2,1.6)}{$\frac{y_0}{\beta_1}$}{$\frac{y_0}{\beta_2}$}
  \node[punto] at (2,1.6) {};
\end{tikzpicture}
\end{document}
